\chapter{Modules and Linear Algebra}\label{chap:modules-linear-algebra}

Vector spaces combine addition with scalar multiplication by numbers from a
field. Modules keep the same formal operations while allowing scalars from a
ring. This modest-looking change includes abelian groups, ideals, quotient
rings, and vector spaces in one language. It also explains why linear algebra
is the first important special case of module theory: over a field one can
divide by every nonzero scalar, whereas a general module can have torsion and
need not possess coordinates.

Throughout \(R\) is a ring with identity, and a module means a unital left
\(R\)-module unless a different convention is stated. For a noncommutative
ring this is important: a right module has its scalar multiplication on the
other side. From \S6.8 onward, \(F\) denotes a field and \(V,W\) denote
finite-dimensional \(F\)-vector spaces.

\section{Modules and Submodules}
The vector-space axioms make sense before one assumes that nonzero scalars
are invertible. That observation gives the correct setting for relations such
as \(n\bar a=0\) in \(\mathbb Z/n\mathbb Z\), which cannot occur in a
nonzero vector space.

\begin{definition}
An \emph{\(R\)-module} is an abelian group \((M,+)\) equipped with a map
\(R\times M\to M\), \((r,m)\mapsto rm\), such that, for \(r,s\in R\) and
\(m,n\in M\),
\[
\begin{aligned}
(r+s)m&=rm+sm, & r(m+n)&=rm+rn,\\
(rs)m&=r(sm), & 1m&=m.
\end{aligned}
\]
A subset \(N\subseteq M\) is a \emph{submodule}, written \(N\leq M\), if
\(m-n\in N\) and \(rm\in N\) whenever \(m,n\in N\) and \(r\in R\).
An \emph{\(R\)-module homomorphism} \(f:M\to P\) preserves addition and
scalar multiplication. Its \emph{kernel} and \emph{image} are
\(\ker f=\{m:f(m)=0\}\) and \(\operatorname{im}f=\{f(m):m\in M\}\).
An isomorphism is a bijective module homomorphism.
\end{definition}

The principal examples should be kept in view. A vector space is a module
over its scalar field; abelian groups are exactly \(\mathbb Z\)-modules;
\(R\) and \(R^n\) are left \(R\)-modules; and every left ideal is a
submodule of the regular module \(R\). If \(R\) is commutative, every ideal
has this role. The module \(\mathbb Z/n\mathbb Z\) is not free over
\(\mathbb Z\), because its nonzero elements are killed by a nonzero scalar.

There is a useful one-line test hidden in the definition. A nonempty subset
\(N\) of a module is a submodule precisely when it is closed under
subtraction and under multiplication by every scalar. The subtraction test
contains both additive closure and additive inverses, just as it did for
subgroups. For instance, \(2\mathbb Z\leq\mathbb Z\) as a
\(\mathbb Z\)-module, while the set of odd integers is not a submodule:
the difference of two odd integers is even. Such examples are elementary,
but they train the reader to distinguish closure conditions from merely
familiar-looking subsets.

\begin{proposition}
For a module homomorphism \(f:M\to P\), both \(\ker f\) and
\(\operatorname{im}f\) are submodules.
\end{proposition}
\begin{proof}
If \(m,n\in\ker f\), then \(f(m-n)=0\), and \(f(rm)=rf(m)=0\). Thus the
kernel is a submodule. If \(f(m),f(n)\in\operatorname{im}f\), then
\(f(m)-f(n)=f(m-n)\) and \(rf(m)=f(rm)\), so the image is one as well.
\end{proof}

\section{Quotient Modules and the Isomorphism Theorems}
Submodules are precisely the pieces that can be collapsed to zero without
destroying scalar multiplication. Quotient modules therefore generalize the
quotient groups and quotient rings developed earlier.

\begin{definition}
If \(N\leq M\), the \emph{quotient module} \(M/N\) is the additive quotient
group with scalar multiplication \(r(m+N)=rm+N\).
\end{definition}
\begin{proof}
If \(m+N=m'+N\), then \(m-m'\in N\), whence \(r(m-m')\in N\). Therefore
\(rm+N=rm'+N\), so the operation is well-defined. Each module axiom follows
from the corresponding axiom in \(M\).
\end{proof}

\begin{theorem}[Module isomorphism theorems]
Let \(f:M\to P\) be linear.
\begin{enumerate}
\item \(M/\ker f\cong\operatorname{im}f\).
\item If \(A,B\leq M\), then \((A+B)/B\cong A/(A\cap B)\).
\item If \(N\leq L\leq M\), then \((M/N)/(L/N)\cong M/L\).
\item Submodules of \(M/N\) correspond bijectively to submodules of \(M\)
containing \(N\), under inverse image and quotient by \(N\).
\end{enumerate}
\end{theorem}
\begin{proof}
For (1), \(m+\ker f\mapsto f(m)\) is well-defined, onto the image, and has
zero kernel. For (2), restrict \(M\to M/B\) to \(A\): its kernel is
\(A\cap B\), and its image is \((A+B)/B\). For (3), map
\((m+N)+(L/N)\) to \(m+L\), checking directly that the map is well-defined
and has zero kernel. Finally, if \(\pi:M\to M/N\) is the quotient map,
then \(Q\mapsto\pi^{-1}(Q)\) and \(K\mapsto K/N\) are inverse operations.
\end{proof}

\begin{theorem}[Universal property of a quotient]
If \(f:M\to P\) is linear and \(N\subseteq\ker f\), there is a unique
linear map \(\bar f:M/N\to P\) with \(f=\bar f\pi\), where \(\pi\) is the
quotient map.
\end{theorem}
\[
\begin{tikzcd}[column sep=large, row sep=large]
M \arrow[r, "\pi"] \arrow[dr, "f"'] & M/N \arrow[d, "\bar f"] \\
& P
\end{tikzcd}
\]
\begin{proof}
Set \(\bar f(m+N)=f(m)\). If \(m-m'\in N\), then \(f(m-m')=0\), so this
does not depend on the representative. It is linear, and surjectivity of
\(\pi\) forces uniqueness.
\end{proof}

Thus the group, ring, and module isomorphism theorems are not unrelated
coincidences: they express one recurring quotient pattern.

\begin{example}[A quotient module]
The quotient \(\mathbb Z/6\mathbb Z\) is a \(\mathbb Z\)-module, and its
submodule generated by \(\bar2\) is \(\{\bar0,\bar2,\bar4\}\). Collapsing
this submodule gives a quotient with two elements,
\[(\mathbb Z/6\mathbb Z)/\langle\bar2\rangle\cong\mathbb Z/2\mathbb Z.
\]
The example is the module version of reducing a group by a normal subgroup;
the point of the module language is that multiplication by every integer also
descends automatically.
\end{example}

\section{Direct Sums and Products}
Products collect arbitrary coordinate choices; direct sums collect finite
linear combinations of contributions from separate modules. The latter is the
algebraic construction needed for bases and for decompositions.

\begin{definition}
For a family \((M_i)_{i\in I}\), the product
\[
\prod_{i\in I}M_i
\]
consists of all tuples \((m_i)\), with coordinatewise operations. The direct
sum
\[
\bigoplus_{i\in I}M_i
\]
is its submodule consisting of
tuples with \emph{finite support}: only finitely many \(m_i\) are nonzero.
For a finite family, product and direct sum have the same elements and are
both denoted
\[
M_1\oplus\cdots\oplus M_n.
\]
\end{definition}

There are canonical injections and projections
\[
\iota_i:M_i\longrightarrow\bigoplus_{i\in I}M_i,
\qquad
\pi_i:\prod_{i\in I}M_i\longrightarrow M_i.
\]
A map
\[
X\longrightarrow\prod_{i\in I}M_i
\]
is
equivalent to a family of maps \(X\to M_i\), obtained by composing with the
projections. Conversely, a family \(M_i\to Y\) determines a unique map
\[
\bigoplus_{i\in I}M_i\longrightarrow Y,
\]
because every element is a finite sum of injected
components. For example,
\(\mathbb Z/2\mathbb Z\oplus\mathbb Z/3\mathbb Z\cong\mathbb Z/6\mathbb Z\).

For an infinite family the finite-support requirement matters. The element
\[
(1,1,1,\ldots)\in\prod_{n\geq1}\mathbb Z
\]
does not belong to
\[
\bigoplus_{n\geq1}\mathbb Z.
\]
A map out of the direct sum can be defined
by adding only finitely many contributions, whereas there is no reason an
infinite sum should make sense in an arbitrary module. This is why direct
sums, rather than products, are the natural home for bases.

\section{Cyclic Modules, Annihilators, Torsion, and Exact Sequences}
The simplest module is generated by one element. Understanding such modules
is useful because a generator has exactly one kind of obstruction: some
scalars may kill it. The annihilator records those scalar relations and turns
cyclic modules into quotients of the ring itself.

\begin{definition}
An \(R\)-module is \emph{cyclic} if \(M=Rm\) for some \(m\in M\). The
\emph{annihilator} of \(m\) is
\[\operatorname{Ann}_R(m)=\{r\in R:rm=0\};\]
the annihilator of \(M\) consists of scalars that kill every element of
\(M\). If \(R\) is a domain, an element \(m\) is \emph{torsion} if
\(rm=0\) for some nonzero \(r\in R\). A module is torsion-free if its only
torsion element is zero.
\end{definition}

For a left module, an annihilator is a left ideal; in the commutative rings
used later it is an ideal in the unqualified sense. Over \(\mathbb Z\),
torsion means finite additive order. Thus \(\mathbb Z^n\) is torsion-free,
whereas every element of \(\mathbb Z/n\mathbb Z\) is torsion.

When \(R\) is a domain, the set \(T(M)\) of torsion elements is a submodule,
called the \emph{torsion submodule}. Indeed, if \(am=0\) and \(bn=0\) with
\(a,b\ne0\), then \(ab(m-n)=0\); and \(a(rm)=0\) for every \(r\in R\).
The domain hypothesis ensures \(ab\ne0\). This construction will make the
free part of a finitely generated PID module intrinsic in Chapter~7.

\begin{theorem}[Cyclic-module classification]
For \(m\in M\), there is an isomorphism
\[Rm\cong R/\operatorname{Ann}_R(m).\]
Consequently every cyclic module over a commutative ring is isomorphic to
\(R/I\) for an ideal \(I\).
\end{theorem}
\begin{proof}
The map \(R\to Rm\), \(r\mapsto rm\), is a surjective module
homomorphism, and its kernel is exactly \(\operatorname{Ann}_R(m)\). The
first isomorphism theorem supplies the stated isomorphism. Conversely,
\(R/I\) is generated by \(1+I\).
\end{proof}

Exact sequences provide a compact language for saying that a module is built
from a submodule and a quotient. They will later organize module
decompositions and further constructions.
\begin{definition}
A sequence \(A\xrightarrow{u}B\xrightarrow{v}C\) is \emph{exact at
\(B\)} if \(\operatorname{im}u=\ker v\). A \emph{short exact sequence}
\[0\longrightarrow A\xrightarrow{i}B\xrightarrow{q}C\longrightarrow0\]
is exact everywhere; equivalently, \(i\) is injective, \(q\) is surjective,
and \(\operatorname{im}i=\ker q\).
\end{definition}

\begin{undefinition}[Cokernel]
For a module homomorphism \(f:M\to N\), its \emph{cokernel} is
\[
\operatorname{coker}(f)=N/\operatorname{im}f.
\]
The canonical map \(N\to\operatorname{coker}(f)\) is the quotient map.
\end{undefinition}

Kernels measure what a map kills; cokernels measure what remains in its target
after the image is collapsed.  Concretely, if \(g:N\to P\) satisfies
\(gf=0\), then \(g\) vanishes on \(\operatorname{im}f\), so there is a
unique map \(\widetilde g:\operatorname{coker}(f)\to P\) with
\[
\begin{tikzcd}[column sep=large,row sep=large]
N\arrow[r]\arrow[dr,"g"']&\operatorname{coker}(f)\arrow[d,"\widetilde g"]\\
&P.
\end{tikzcd}
\]

The quotient sequence \(0\to N\to M\to M/N\to0\) is the basic example.
It is useful to distinguish equality of short exact sequences from the
weaker fact that their three terms happen to be isomorphic.

\begin{definition}
An \emph{isomorphism of short exact sequences} is a commutative diagram
\[
\begin{tikzcd}[column sep=large]
0\arrow[r]&A\arrow[r,"i"]\arrow[d,"\alpha", "\sim"']&
B\arrow[r,"q"]\arrow[d,"\beta", "\sim"']&
C\arrow[r]\arrow[d,"\gamma", "\sim"']&0\\
0\arrow[r]&A'\arrow[r,"i'"]&B'\arrow[r,"q'"]&C'\arrow[r]&0
\end{tikzcd}
\]
in which the three vertical maps are module isomorphisms.
\end{definition}

\begin{proposition}[The canonical split sequence]\label{prop:canonical-split-sequence}
For modules \(A\) and \(C\), the sequence
\[
0\longrightarrow A\xrightarrow{a\mapsto(a,0)}A\oplus C
\xrightarrow{(a,c)\mapsto c}C\longrightarrow0
\]
is short exact.
\end{proposition}
\begin{proof}
The first map is injective and the second is surjective.  Its kernel consists
precisely of the pairs \((a,0)\), which are exactly the image of the first
map.  Thus exactness says, quite literally, that \(A\oplus C\) contains
\(A\) as its first coordinate and has quotient \(C\).
\end{proof}

\section{Splitting and Diagram Chases}
The preceding observation asks when a short exact sequence contains no
extension information beyond its two ends.  A map \(s:C\to B\) with
\(qs=1_C\) is a \emph{section}; a map \(r:B\to A\) with \(ri=1_A\) is a
\emph{retraction}.

\begin{theorem}[Splitting Lemma]\label{thm:splitting-lemma}
For a short exact sequence
\[
0\longrightarrow A\xrightarrow{i}B\xrightarrow{q}C\longrightarrow0,
\]
the following are equivalent: \(q\) has a section; \(i\) has a retraction;
and the given short exact sequence is isomorphic, as a short exact sequence,
to the canonical split sequence of
Proposition~\ref{prop:canonical-split-sequence}.  Concretely, the end maps
are \(1_A\) and \(1_C\), and the middle isomorphism
\(\Phi:A\oplus C\to B\) satisfies
\[
\Phi(a,0)=i(a),
\qquad
q\Phi(a,c)=c.
\]
\end{theorem}
\begin{proof}
A section gives the map
\[
\Phi:A\oplus C\longrightarrow B,\qquad
\Phi(a,c)=i(a)+s(c).
\]
It is onto: for \(b\in B\), the element \(b-sq(b)\) lies in \(\ker q\), hence
is \(i(a)\) for some \(a\).  If \(\Phi(a,c)=0\), applying \(q\) gives
\(c=0\), and then injectivity of \(i\) gives \(a=0\).  Thus \(\Phi\) is an
isomorphism, and it gives an isomorphism to the canonical sequence in
Proposition~\ref{prop:canonical-split-sequence}.  The map
\(r=\operatorname{pr}_A\Phi^{-1}\) is a retraction of \(i\).

Conversely, suppose that \(ri=1_A\).  For \(c\in C\), choose \(b\in B\) with
\(q(b)=c\), and put
\[
s(c)=b-i(r(b)).
\]
If \(b'=b+i(a)\) is another lift, then
\[
b'-i(r(b'))=b+i(a)-i(r(b)+a)=b-i(r(b)),
\]
so the formula is independent of the lift.  It is linear, and
\(q(s(c))=c\).  Hence it is a section.
\end{proof}

Commutative diagrams organize a \emph{diagram chase}: start with an element,
use exactness to replace a kernel element by a preimage, and use commutativity
to transport the resulting equality vertically.  The Five Lemma is the basic
instance.

\begin{theorem}[Five Lemma]\label{thm:five-lemma}
In a commutative diagram with exact rows
\[
\begin{tikzcd}[column sep=small]
A_1\arrow[r]\arrow[d,"f_1"]&A_2\arrow[r]\arrow[d,"f_2"]&A_3\arrow[r]\arrow[d,"f_3"]&A_4\arrow[r]\arrow[d,"f_4"]&A_5\arrow[d,"f_5"]\\
B_1\arrow[r]&B_2\arrow[r]&B_3\arrow[r]&B_4\arrow[r]&B_5
\end{tikzcd}
\]
if \(f_1\) is surjective, \(f_2,f_4\) are isomorphisms, and \(f_5\) is
injective, then \(f_3\) is an isomorphism.
\end{theorem}
\begin{proof}
Write the horizontal maps as \(d_i:A_i\to A_{i+1}\) and
\(e_i:B_i\to B_{i+1}\).  Suppose first that \(f_3(x)=0\).  Then
\(f_4d_3(x)=e_3f_3(x)=0\).  Since \(f_4\) is injective,
\(d_3(x)=0\), so \(x=d_2(y)\).  As
\[
0=f_3d_2(y)=e_2f_2(y),
\]
there is \(z\in B_1\) with \(e_1(z)=f_2(y)\).  Choose \(w\in A_1\) with
\(f_1(w)=z\).  Then
\[
f_2(y-d_1(w))=0.
\]
Injectivity of \(f_2\) gives \(y=d_1(w)\), hence \(x=d_2d_1(w)=0\).
This proves the \emph{injectivity} of \(f_3\).

\emph{Surjectivity.} Now take \(b\in B_3\).  Since \(e_3(b)\in B_4\), choose
\(a_4\in A_4\) with \(f_4(a_4)=e_3(b)\).  Exactness gives
\(d_4(a_4)=0\), because \(f_5d_4(a_4)=e_4f_4(a_4)=0\) and \(f_5\) is
injective.  Choose \(a_3\in A_3\) with \(d_3(a_3)=a_4\).  Then
\[
e_3(b-f_3(a_3))=0,
\]
so \(b-f_3(a_3)=e_2(b_2)\) for some \(b_2\in B_2\).  Choose
\(a_2\in A_2\) with \(f_2(a_2)=b_2\).  Replacing \(a_3\) by
\(a_3+d_2(a_2)\) gives \(f_3(a_3)=b\).  This proves the \emph{surjectivity}
of \(f_3\).

\emph{Isomorphism.} A module homomorphism is an isomorphism precisely when it
is both injective and surjective.  The preceding two parts prove the claim.
\end{proof}

\begin{theorem}[Snake Lemma]\label{thm:snake-lemma}
In a commutative diagram with exact rows
\[
\begin{tikzcd}[column sep=large]
0\arrow[r]&A'\arrow[r,"i'"]\arrow[d,"f_A"]&
B'\arrow[r,"q'"]\arrow[d,"f_B"]&
C'\arrow[r]\arrow[d,"f_C"]&0\\
0\arrow[r]&A\arrow[r,"i"]&
B\arrow[r,"q"]&
C\arrow[r]&0,
\end{tikzcd}
\]
there is an exact sequence
\[
\begin{gathered}
0\longrightarrow\ker f_A\longrightarrow\ker f_B
\longrightarrow\ker f_C\xrightarrow{\delta}\operatorname{coker}f_A\\
\longrightarrow\operatorname{coker}f_B\longrightarrow
\operatorname{coker}f_C\longrightarrow0.
\end{gathered}
\]
\end{theorem}
\begin{proof}
For \(c'\in\ker f_C\), choose \(b'\in B'\) with \(q'(b')=c'\).  Then
\[
q(f_B(b'))=f_C(q'(b'))=0,
\]
so \(f_B(b')=i(a)\) for a unique \(a\in A\).  Define
\[
\delta(c')=a+\operatorname{im}f_A.
\]
If \(b'\) is replaced by \(b'+i'(a')\), then \(a\) is replaced by
\(a+f_A(a')\); hence \(\delta\) is well-defined.  Choosing lifts of two
elements and adding them also proves that \(\delta\) is linear.

The maps before and after \(\delta\) are the maps induced by the vertical
arrows.  The first one is injective because \(i'\) is injective.  Exactness
at the first two kernels follows immediately by chasing the two squares and
using injectivity of \(i\).  If \(c'=q'(b')\) comes from
\(\ker f_B\), then \(f_B(b')=0\), so \(\delta(c')=0\).  Conversely,
\(\delta(c')=0\) means that \(a=f_A(a')\); then
\(b'-i'(a')\) lies in \(\ker f_B\) and still maps to \(c'\).  This proves
exactness at \(\ker f_C\).

If \(a+\operatorname{im}f_A\) maps to zero in \(\operatorname{coker}f_B\),
then \(i(a)=f_B(b')\) for some \(b'\).  Hence
\(f_C(q'(b'))=0\), and the construction gives
\(\delta(q'(b'))=a+\operatorname{im}f_A\).  This is exactness at
\(\operatorname{coker}f_A\).  Finally, exactness at
\(\operatorname{coker}f_B\) and \(\operatorname{coker}f_C\) follows by the
same two-square chase: a class \(b+\operatorname{im}f_B\) maps to zero
precisely when \(q(b)=f_C(c')\), and a lift of \(c'\) corrects \(b\) to an
element of \(i(A)\).  The last induced map is surjective because \(q\) is
surjective.  This also proves the asserted endpoint exactness.
\end{proof}

The construction of \(\delta\) is the model diagram chase: choose a lift,
measure the obstruction in the next kernel, and record it modulo the
ambiguity forced by the first vertical map.  It is natural with respect to
maps of such diagrams, a fact used later in homological algebra.

\section{Free Modules and Module Bases}
Coordinates are useful only when expansions are unique. Free modules record
that condition over a general ring; vector spaces are the especially well
behaved free modules over fields.

\begin{definition}
A subset \(S\subseteq M\) \emph{spans} \(M\) if every element is a finite
linear combination of members of \(S\). It is \emph{linearly independent}
if every relation
\[
\sum_i r_is_i=0,
\]
with distinct \(s_i\in S\), has all
coefficients zero. A \emph{basis} is an independent spanning set, and a
module with a basis is \emph{free}. The free module on a set \(X\) is the
module \(R^{(X)}\) of finitely supported functions \(X\to R\), with basis
\(\{e_x:x\in X\}\).
\end{definition}

\begin{theorem}[Universal property of a free module]
For every function \(u:X\to M\), there is a unique linear map
\(R^{(X)}\to M\) taking \(e_x\) to \(u(x)\).
\end{theorem}
\begin{proof}
Every element of \(R^{(X)}\) has a unique finite expression
\[
\sum_{x\in X}r_xe_x.
\]
Send it to
\[
\sum_{x\in X}r_xu(x).
\]
This is linear, and any
linear map with the required values has this formula.
\end{proof}

Not every module is free: \(\mathbb Z/n\mathbb Z\) has no \(\mathbb Z\)-basis
for \(n>1\). Indeed, if \(\bar a\) belonged to a \(\mathbb Z\)-basis, then
\(n\bar a=0\) would be a nontrivial relation. This is the first warning that
the familiar basis theory of vector spaces cannot simply be imported into
module theory. In particular, a submodule of a free module need not be free
over an arbitrary ring, and the apparent number of generators of \(R^n\)
requires justification before it is called a rank.

\section{Projective and Injective Modules}
Free modules make lifting possible because a map out of a free module is
controlled by its values on a basis.  Projective modules are precisely the
modules that retain this lifting property, even when they need not themselves
have a basis.

\begin{definition}
A module \(P\) is \emph{projective} if, for every surjection
\(\pi:M\twoheadrightarrow N\) and every map \(f:P\to N\), there is a map
\(\widetilde f:P\to M\) such that \(\pi\widetilde f=f\).  A module \(E\) is
\emph{injective} if, for every injection \(j:A\hookrightarrow B\) and every
map \(f:A\to E\), there is a map \(\widetilde f:B\to E\) such that
\(\widetilde f j=f\).
\end{definition}
\[
\begin{tikzcd}[column sep=large,row sep=large]
&P\arrow[dl,dashed,"\widetilde f"']\arrow[d,"f"]\\
M\arrow[r,two heads,"\pi"']&N\arrow[r]&0
\end{tikzcd}
\qquad
\begin{tikzcd}[column sep=large,row sep=large]
&E\\
0\arrow[r]&A\arrow[u,"f"]\arrow[r,hook,"j"']&B\arrow[ul,dashed,"\widetilde f"']
\end{tikzcd}
\]
These matching triangles display the dual geometry of the two definitions:
projectivity is a lifting problem, while injectivity is an extension problem.

\begin{definition}
For left \(R\)-modules \(M,N\), write
\[
\operatorname{Hom}_R(M,N)
\]
for the set of \(R\)-linear maps \(M\to N\).  It is an abelian group under
pointwise addition:
\[
(f+g)(m)=f(m)+g(m).
\]
If \(R\) is commutative, it is also an \(R\)-module by
\[
(rf)(m)=r f(m).
\]
\end{definition}
For a noncommutative ring the displayed scalar formula need not define a
left-module map, since \(r f(sm)\) need not equal \(s r f(m)\).  Accordingly,
the exact sequences below are always sequences of abelian groups, and only
under commutativity are they automatically sequences of \(R\)-modules.

\begin{proposition}[Hom detects exactness]\label{prop:hom-exactness}
Suppose
\[
0\longrightarrow A\xrightarrow{i}B\xrightarrow{q}C\longrightarrow0
\]
satisfies \(qi=0\), without yet assuming exactness.  Define
\[
\begin{aligned}
i_*:\operatorname{Hom}_R(M,A)&\longrightarrow\operatorname{Hom}_R(M,B),
&f&\longmapsto i\circ f,\\
q_*:\operatorname{Hom}_R(M,B)&\longrightarrow\operatorname{Hom}_R(M,C),
&g&\longmapsto q\circ g,
\end{aligned}
\]
\[
\begin{aligned}
q^*:\operatorname{Hom}_R(C,N)&\longrightarrow\operatorname{Hom}_R(B,N),
&h&\longmapsto h\circ q,\\
i^*:\operatorname{Hom}_R(B,N)&\longrightarrow\operatorname{Hom}_R(A,N),
&k&\longmapsto k\circ i.
\end{aligned}
\]
Then
\[
0\longrightarrow\operatorname{Hom}_R(M,A)\xrightarrow{i_*}
\operatorname{Hom}_R(M,B)\xrightarrow{q_*}\operatorname{Hom}_R(M,C)
\]
is exact for every \(R\)-module \(M\) if and only if
\[
0\longrightarrow A\xrightarrow{i}B\xrightarrow{q}C
\]
is exact.  Dually,
\[
0\longrightarrow\operatorname{Hom}_R(C,N)\xrightarrow{q^*}
\operatorname{Hom}_R(B,N)\xrightarrow{i^*}\operatorname{Hom}_R(A,N)
\]
is exact for every \(R\)-module \(N\) if and only if
\[
A\xrightarrow{i}B\xrightarrow{q}C\longrightarrow0
\]
is exact.  Thus the original sequence is short exact precisely when both
displayed Hom tests are exact for all \(M\) and \(N\).
\end{proposition}
\begin{proof}
Assume first that \(0\to A\to B\to C\) is exact.  A map \(M\to B\) has
zero composite with \(q\) exactly when its image lies in
\(\ker q=\operatorname{im}i\), and it then factors uniquely through \(A\);
this proves exactness of the first Hom sequence.  Conversely, suppose this
Hom sequence is exact for every \(M\).  With \(K=\ker i\), the inclusion
\(K\to A\) has zero image under \(i_*\), hence is zero; so \(i\) is
injective.  With \(K=\ker q\), the inclusion \(K\to B\) has zero image
under \(q_*\), hence factors through \(i\).  Thus
\(\ker q\subseteq\operatorname{im}i\), and the reverse inclusion follows
from \(qi=0\).

Now assume that \(A\to B\to C\to0\) is exact.  A map \(B\to N\) vanishes on
\(\operatorname{im}i=\ker q\) exactly when it factors uniquely through
\(C\), proving exactness of the second Hom sequence.  Conversely, suppose
that sequence is exact for every \(N\).  Taking
\[
N=C/\operatorname{im}q
\]
and the quotient map \(C\to N\), injectivity of \(q^*\) forces that quotient
map to be zero; hence \(q\) is surjective.  Taking
\[
N=B/\operatorname{im}i
\]
and the quotient map \(B\to N\), exactness gives a factorization through
\(q\).  It vanishes on \(\ker q\), so
\(\ker q\subseteq\operatorname{im}i\); the opposite inclusion again follows
from \(qi=0\).
\end{proof}

For a short exact sequence, both Hom sequences are exact at their initial
terms.  This familiar property is often called left exactness of Hom.
Exactness through the final Hom term is strictly stronger: with \(M=P\) it
is the lifting property defining projectivity, and with \(N=E\) it is the
extension property defining injectivity.

\begin{theorem}[Projective-module characterizations]\label{thm:projective-characterizations}
For an \(R\)-module \(P\), the following are equivalent:
\begin{enumerate}
\item \(P\) is projective;
\item every short exact sequence ending in \(P\) splits;
\item \(P\) is isomorphic to a direct summand of a free module;
\item for every short exact sequence, the first sequence in
Proposition~\ref{prop:hom-exactness}, with \(M=P\), is exact at its final
term.
\end{enumerate}
In particular, every free module is projective.
\end{theorem}
\begin{proof}
If \(P\) is projective and
\[
0\longrightarrow A\longrightarrow B\xrightarrow{q}P\longrightarrow0
\]
is exact, apply the lifting property to \(1_P:P\to P\).  The resulting
section of \(q\) splits the sequence by the Splitting Lemma.

Suppose every such sequence splits.  Choose a generating set \(X\) of \(P\).
The universal property of \(R^{(X)}\) gives a surjection
\(R^{(X)}\twoheadrightarrow P\).  Its kernel produces a short exact
sequence ending in \(P\), so it splits.  Thus
\[
R^{(X)}\cong\ker(R^{(X)}\to P)\oplus P,
\]
and \(P\) is a direct summand of a free module.

Finally, let \(P\) be a direct summand of a free module \(F\), with
inclusion \(u:P\to F\) and retraction \(r:F\to P\).  Given
\(\pi:M\twoheadrightarrow N\) and \(f:P\to N\), choose a basis of \(F\).
For each basis vector \(e_x\), choose a lift in \(M\) of
\(f(r(e_x))\).  The universal property of \(F\) gives
\(g:F\to M\) with \(\pi g=fr\).  Then \(\widetilde f=gu\) satisfies
\(\pi\widetilde f=f\).  Thus \(P\) is projective.  Taking \(P=F\) proves
the final assertion.  Finally, Proposition~\ref{prop:hom-exactness} says
exactly that the final Hom map is surjective precisely when every map
\(P\to C\) lifts through \(B\).  This proves the equivalence of (1) and
(4), and distinguishes full exactness here from the left exactness enjoyed
by every test module.
\end{proof}

\noindent\textbf{Example.} For \(n>1\), the \(\mathbb Z\)-module
\(\mathbb Z/n\mathbb Z\) is not projective.  If
\[
0\longrightarrow\mathbb Z\xrightarrow{\times n}\mathbb Z
\longrightarrow\mathbb Z/n\mathbb Z\longrightarrow0
\]
split, then \(\mathbb Z\) would contain a subgroup isomorphic to
\(\mathbb Z/n\mathbb Z\).  This is impossible because \(\mathbb Z\) has no
nonzero torsion.

\begin{theorem}[Injective-module characterization]\label{thm:injective-characterization}
An \(R\)-module \(E\) is injective if and only if every short exact sequence
\[
0\longrightarrow E\longrightarrow B\longrightarrow C\longrightarrow0
\]
splits.
\; Equivalently, for every short exact sequence, the second sequence in
Proposition~\ref{prop:hom-exactness}, with \(N=E\), is exact at its final
term.
\end{theorem}
\begin{proof}
If \(E\) is injective, apply its extension property to the identity map of
\(E\) and the first injection in the sequence.  The resulting retraction
splits the sequence.  Conversely, let \(j:A\hookrightarrow B\) and
\(f:A\to E\) be given.  Form
\[
S=\{(f(a),-j(a)):a\in A\}\subseteq E\oplus B,
\qquad
Q=(E\oplus B)/S.
\]
The set \(S\) is a submodule: it is the image of the linear map
\(A\to E\oplus B\), \(a\mapsto(f(a),-j(a))\).  The map
\(E\to Q\), \(e\mapsto(e,0)+S\), is injective: if
\[
(e,0)=(f(a),-j(a))\in S,
\]
then \(j(a)=0\), so \(a=0\) and \(e=0\).  The map
\[
Q\longrightarrow B/j(A),
\qquad
(e,b)+S\longmapsto b+j(A)
\]
is well-defined, since adding \((f(a),-j(a))\) changes \(b\) by an element
of \(j(A)\).  Its kernel is the image of \(E\): if \(b=j(a)\), then
\[
(e,b)+S=(e+f(a),0)+S,
\]
and the reverse inclusion is immediate.  It is onto, so the First
Isomorphism Theorem identifies the quotient of \(Q\) by the image of \(E\)
with \(B/j(A)\).  By hypothesis the resulting short exact sequence splits,
giving a retraction \(Q\to E\).  Composing
\(B\to Q\), \(b\mapsto(0,b)+S\), with that retraction gives an extension of
\(f\).  Hence \(E\) is injective.  The final Hom characterization follows
from Proposition~\ref{prop:hom-exactness}: surjectivity of its final map is
exactly the assertion that maps \(A\to E\) extend to \(B\).
\end{proof}

\begin{proposition}[Enough projectives]\label{prop:enough-projectives}
Every \(R\)-module is a quotient of a free, hence projective, \(R\)-module.
\end{proposition}
\begin{proof}
Let \(X\) be the underlying set of \(M\).  The free module
\[
R^{(X)}
\]
has basis \(\{e_x:x\in X\}\).  The rule \(e_x\mapsto x\) extends uniquely to
a surjective map \(R^{(X)}\to M\).  The final assertion follows from
Theorem~\ref{thm:projective-characterizations}.
\end{proof}

\begin{theorem}[Baer's criterion]\label{thm:baer-criterion}
An \(R\)-module \(E\) is injective if and only if every \(R\)-linear map
\[
f:J\longrightarrow E
\]
from a left ideal \(J\) of \(R\) extends to a map \(R\to E\).
\end{theorem}
\begin{proof}
The forward implication applies injectivity to \(J\hookrightarrow R\).
Conversely, given \(A\subseteq B\) and \(f:A\to E\), order extensions
\[
(C,g),\qquad A\subseteq C\subseteq B,\qquad g|_A=f,
\]
by extension.  Unions of chains give extensions, so Zorn's Lemma gives a
maximal pair \((C,g)\).  If \(b\notin C\), then
\[
J=\{r\in R:rb\in C\}
\]
is a left ideal.  The map \(r\mapsto g(rb)\) extends to \(r\mapsto re\) on
\(R\).  Hence
\[
g'(c+rb)=g(c)+re
\]
is well-defined: if \(c+rb=c'+r'b\), then
\[
(r-r')b=c'-c\in C,
\]
so \(r-r'\in J\).  The chosen extension agrees with
\(s\mapsto g(sb)\) on \(J\), hence
\[
(r-r')e=g((r-r')b)=g(c'-c).
\]
Rearranging gives \(g(c)+re=g(c')+r'e\).  Thus \(g'\) extends \(g\) to
\(C+Rb\), contradicting maximality.  Therefore \(C=B\).
\end{proof}

\begin{remark}
The standard converse proof of Baer's criterion uses Zorn's Lemma; its
forward implication does not.  Over Noetherian rings, left ideals are
finitely generated, so many applications reduce ideal-level arguments to
finite-generation arguments.  The precise foundational strength of the full
theorem is not pursued here.  Looking ahead, every module embeds in an
injective module; standard proofs use additional constructions, often
injective abelian groups such as \(\mathbb Q/\mathbb Z\) and suitable Hom
modules.
\end{remark}

\begin{proposition}[Injective abelian groups]\label{prop:injective-divisible}
An abelian group is injective as a \(\mathbb Z\)-module if and only if it is
\emph{divisible}: multiplication by every nonzero integer is surjective.
\end{proposition}
\begin{proof}
If \(E\) is injective, extend
\[
n\mathbb Z\longrightarrow E,\qquad nk\longmapsto ka
\]
along \(n\mathbb Z\hookrightarrow\mathbb Z\); the image of \(1\) is an
\(n\)-th root of \(a\).  Conversely, every ideal of \(\mathbb Z\) is
\(n\mathbb Z\), and divisibility extends any map from it by sending \(1\) to
an \(n\)-th root of the image of \(n\).  Apply Baer's criterion.
\end{proof}

\begin{example}
\(\mathbb Q\), \(\mathbb Q/\mathbb Z\), and
\[
\mathbb Q^r
\]
are divisible, hence injective \(\mathbb Z\)-modules; \(\mathbb Z\) is not.
\end{example}

\begin{proposition}[Enough injectives over \(\mathbb Z\)]
Every abelian group embeds in a divisible abelian group.  Equivalently, every
\(\mathbb Z\)-module embeds in an injective \(\mathbb Z\)-module.
\end{proposition}
\begin{proof}
The group \(\mathbb Q/\mathbb Z\) is divisible, hence injective by
Proposition~\ref{prop:injective-divisible}.  For every nonzero \(a\in A\),
there is a homomorphism
\[
\chi_a:A\longrightarrow\mathbb Q/\mathbb Z
\]
with \(\chi_a(a)\ne0\).  Indeed, first define such a map on the cyclic
subgroup generated by \(a\): send a generator of a finite cyclic group of
order \(n\) to \(1/n+\mathbb Z\), and send an infinite cyclic generator to
\(1/2+\mathbb Z\).  Injectivity of \(\mathbb Q/\mathbb Z\) extends this map
to \(A\).

Now form the evaluation map
\[
A\longrightarrow
\prod_{\chi\in\operatorname{Hom}_{\mathbb Z}(A,\mathbb Q/\mathbb Z)}
\mathbb Q/\mathbb Z,
\qquad
a\longmapsto\bigl(\chi(a)\bigr)_\chi.
\]
It is injective because the coordinate \(\chi_a\) detects every nonzero
\(a\).  The product is divisible coordinatewise: an \(n\)-th root is chosen
in each coordinate.  Proposition~\ref{prop:injective-divisible} then makes
it injective as a \(\mathbb Z\)-module.
\end{proof}

\begin{remark}
The implication injective \(\Rightarrow\) divisible is elementary.  The
converse, and hence the equivalence used above, is established here through
Baer's criterion and its Zorn-lemma argument.  Thus the presentation relies
on the usual choice principles.
\end{remark}

\section{Vector Spaces, Span, and Finite Bases}
We now specialize to the case in which the scalar ring is a field \(F\).
An \(F\)-vector space is simply an \(F\)-module, so the notions of span,
linear independence, and basis from Sections~6.1 and~6.4 apply without
change. The additional strength of vector-space theory comes from the fact
that every nonzero scalar is invertible.

\begin{definition}
An \emph{\(F\)-vector space} is an \(F\)-module, and a \emph{subspace} is an
\(F\)-submodule.
\end{definition}

\begin{theorem}[Basis existence]
Every vector space has a basis.
\end{theorem}
\begin{proof}
Let \(\mathcal P\) be the set of linearly independent subsets of \(V\),
ordered by inclusion. It is nonempty because it contains the empty set. If
\(\mathcal C\) is a chain in \(\mathcal P\), then
\[
U=\bigcup_{S\in\mathcal C}S
\]
is linearly independent: a finite relation among distinct members of \(U\)
involves only finitely many members of the chain, hence all of them belong to
one member of \(\mathcal C\), where their coefficients vanish. Thus every
chain has an upper bound in \(\mathcal P\). Zorn's lemma gives a maximal
linearly independent set \(B\).

If \(v\notin\operatorname{span}(B)\), then \(B\cup\{v\}\) is linearly
independent: in a relation involving \(v\), a nonzero coefficient of \(v\)
would solve for \(v\) as a linear combination of members of \(B\), while a
zero coefficient leaves a relation in \(B\). This contradicts maximality.
Hence \(B\) spans \(V\), and therefore is a basis.
\end{proof}

\begin{remark}
It applies also to infinite-dimensional spaces. Its proof uses Zorn's lemma
and hence the usual Choice principles. The finite
replacement argument below is elementary and uses no such principle.
\end{remark}

\begin{theorem}[Replacement theorem]
Let
\[
G=\{w_1,\ldots,w_n\}
\]
span \(V\), and let
\[
L=\{v_1,\ldots,v_m\}
\]
be linearly independent. Then \(m\leq n\), and exactly \(n-m\) vectors from
\(G\) can be chosen so that, together with \(v_1,\ldots,v_m\), they span
\(V\). Equivalently, one may replace \(m\) vectors of the original spanning
list by \(v_1,\ldots,v_m\) and retain a spanning list of size \(n\).
\end{theorem}
\begin{proof}
We argue by induction on \(m\). For \(m=0\), the original list spans. Assume
the result for \(m-1\). The induction hypothesis gives a spanning list
\[
\{v_1,\ldots,v_{m-1},w_{i_1},\ldots,w_{i_{n-m+1}}\}.
\]
Write
\[
v_m=a_1v_1+\cdots+a_{m-1}v_{m-1}
+b_1w_{i_1}+\cdots+b_{n-m+1}w_{i_{n-m+1}}.
\]
Not every \(b_j\) is zero, for otherwise \(v_m\) belongs to the span of
\(v_1,\ldots,v_{m-1}\). Choose \(b_k\ne0\). Since \(F\) is a field,
\[
w_{i_k}
=b_k^{-1}\left(
v_m-\sum_{\ell=1}^{m-1}a_\ell v_\ell
-\sum_{j\ne k}b_jw_{i_j}
\right).
\]
Replacing \(w_{i_k}\) by \(v_m\) therefore preserves spanning. There is an
unreplaced \(w_{i_k}\) at every stage, so necessarily \(m\leq n\).
\end{proof}

\begin{corollary}[Extension and extraction]
If \(V\) has a finite spanning set, every finite linearly independent subset
extends to a basis, and every finite spanning set contains a basis.
\end{corollary}
\begin{proof}
Apply the replacement theorem to an independent list and a fixed finite
spanning list. Its retained members, together with the independent list, form
a finite spanning list; discard a retained member whenever it lies in the
span of the others. Conversely, applying the same deletion process to the
original spanning list leaves an independent spanning list.
\end{proof}

\begin{corollary}[Uniqueness of finite basis size]
Any two finite bases of a vector space have the same cardinality.
\end{corollary}
\begin{proof}
Apply the replacement theorem to either basis as an independent set and the
other as a spanning set. The two resulting inequalities are opposite.
\end{proof}

\begin{definition}
A vector space is \emph{finite-dimensional} if it has a finite basis. Its
\emph{dimension}, denoted by \(\dim_FV\), is the common cardinality of its
finite bases.
\end{definition}

Thus a basis expansion exists because a basis spans and is unique because a
basis is independent. An \emph{ordered basis} is a basis displayed as an
ordered finite list; this extra order becomes essential when coordinates are
introduced in Section~6.11.

\begin{proposition}[Subspaces of finite-dimensional spaces]
If \(V\) is finite-dimensional and \(W\leq V\), then \(W\) has a finite
basis and \(\dim_F W\leq\dim_FV\).
\end{proposition}
\begin{proof}
Begin with the empty independent list in \(W\), and adjoin an element of
\(W\) not in its span whenever one exists. The replacement theorem, applied
to the resulting independent list and a fixed basis of \(V\), shows that at
most \(\dim_FV\) choices can be made. When the process stops its list spans
\(W\), so it is a basis. The same theorem gives the inequality.
\end{proof}

\begin{corollary}[IBN for fields]
If
\[
F^m\cong F^n
\]
as \(F\)-vector spaces, then \(m=n\).
\end{corollary}
\begin{proof}
An isomorphism carries the standard basis of \(F^m\) to a basis of \(F^n\).
Uniqueness of finite basis size gives \(m=n\).
\end{proof}

\section{Invariant Basis Number and Finite Free Rank}
The preceding corollary proves IBN for fields. A finite free module over a
general ring needs a separate argument. This distinction is essential in
Chapter~7, where the free part of a module over a PID must have a uniquely
determined rank.

\begin{definition}
A ring \(R\) has \emph{invariant basis number} (IBN) if
\(R^m\cong R^n\) for finite \(m,n\) implies \(m=n\). When \(R\) has IBN,
the \emph{rank} of a finite free \(R\)-module is the integer \(n\) for which
it is isomorphic to \(R^n\).
\end{definition}

\begin{theorem}[Commutative rings have IBN]
Every nonzero commutative ring with identity has invariant basis number.
\end{theorem}
\begin{proof}
Suppose \(R^m\cong R^n\). Since \(R\) is nonzero, \((0)\) is proper; by the
\hyperref[thm:maximal-ideal-existence]{maximal-ideal theorem}, choose a
maximal ideal \(\mathfrak m\) of \(R\). An isomorphism \(f:R^m\to R^n\) carries
\(\mathfrak mR^m\) onto \(\mathfrak mR^n\), since
\(f(rx)=rf(x)\), and the same assertion for \(f^{-1}\) proves equality.
It consequently induces an isomorphism
\[
R^m/\mathfrak mR^m\ \cong\ R^n/\mathfrak mR^n.
\]
Coordinatewise reduction identifies these quotient modules with
\((R/\mathfrak m)^m\) and \((R/\mathfrak m)^n\). The quotient is a field
because \(\mathfrak m\) is maximal, so IBN for fields gives \(m=n\).
\end{proof}

The nonzero hypothesis is essential: for the zero ring,
\[
R^m=R^n=0
\]
for all finite \(m,n\), so invariant basis number fails.  Arbitrary
noncommutative rings need not have IBN, so no unqualified rank is attached to
their finite free modules.

\section{Linear Transformations and Dimension}
The module homomorphisms of vector spaces are the familiar linear
transformations. Finite dimension turns the abstract kernel--image picture
into a numerical conservation law.

\begin{definition}
A \emph{linear transformation} \(T:V\to W\) is an \(F\)-module
homomorphism: \(T(av+bw)=aT(v)+bT(w)\). Its kernel and image are the
submodules already defined in \S6.1, now called respectively its
\emph{kernel} and \emph{range}. It is injective, surjective, bijective, or
an isomorphism in the same sense as a module map. An endomorphism is a map
\(V\to V\), and an automorphism is an invertible endomorphism. The
endomorphisms form the ring \(\mathcal{L}_F(V)\); when the field is clear,
we write \(\mathcal{L}(V)\). Once
\(V\) is finite-dimensional, its \emph{nullity} and \emph{rank} are
\(\dim\ker T\) and \(\dim\operatorname{im}T\).
\end{definition}

The kernel and image are subspaces by Proposition~6.1, and the preceding
subspace theorem supplies their finite bases. Thus the numerical words in the
definition are legitimate before rank--nullity is invoked.

\begin{theorem}[Rank--nullity]
If \(T:V\to W\) is linear and \(V\) is finite-dimensional, then
\[\dim V=\dim\ker T+\dim\operatorname{im}T.\]
\end{theorem}
\begin{proof}
Choose a basis \(k_1,\ldots,k_a\) of \(\ker T\). Extend it to the basis
\[
k_1,\ldots,k_a,v_1,\ldots,v_b
\]
of \(V\). The vectors \(Tv_1,\ldots,Tv_b\)
span the image: applying \(T\) removes the \(k_i\)-part of any coordinate
expression. If
\[
\sum_i c_iTv_i=0,
\]
then
\[
\sum_i c_iv_i\in\ker T,
\]
and
uniqueness of the displayed basis expansion gives every \(c_i=0\). Hence
they form a basis of the image, so \(\dim V=a+b\).
\end{proof}

Thus a linear map between equal finite dimensions is injective exactly when
it is surjective, and hence exactly when it is an isomorphism: rank--nullity
shows that one of kernel and image has the necessary dimension precisely when
the other does.

\section{Matrices and Change of Basis}\label{sec:matrices-change-basis}
Once ordered bases are chosen, a linear transformation becomes a table of
coordinates. The matrix is a useful representative, but it is not the
transformation itself: changing coordinates changes its matrix.

Let \(\beta=(v_1,\ldots,v_n)\) be an ordered basis. Every \(v\in V\) has
one and only one expansion \(v=x_1v_1+\cdots+x_nv_n\); define its
\emph{coordinate column} by
\[
[v]_\beta=\begin{pmatrix}x_1\\x_2\\\vdots\\x_n\end{pmatrix}.
\]
The order matters: exchanging two basis vectors generally exchanges two
coordinates.

\begin{definition}
Let \(R\) be a ring. An \emph{\(m\times n\) matrix over \(R\)} is a
rectangular array
\[
A=(a_{ij})=
\begin{pmatrix}
a_{11}&a_{12}&\cdots&a_{1n}\\
a_{21}&a_{22}&\cdots&a_{2n}\\
\vdots&\vdots&\ddots&\vdots\\
a_{m1}&a_{m2}&\cdots&a_{mn}
\end{pmatrix},
\qquad
a_{ij}\in R.
\]
The set of all such matrices is \(M_{m\times n}(R)\), and
\[
M_n(R)=M_{n\times n}(R).
\]
Matrices are equal precisely when their corresponding entries are equal. For
matrices of the same size, define
\[
A+B=(a_{ij}+b_{ij}).
\]
If \(R\) is commutative, scalar multiplication is defined entrywise by
\[
rA=(ra_{ij}),
\]
making \(M_{m\times n}(R)\) an \(R\)-module.
\end{definition}

Matrices are arrays of scalars; they encode linear maps only after ordered
bases have been chosen. If
\[
A=(a_{ij})\in M_{m\times n}(R),
\]
its \emph{transpose} is the matrix
\[
A^t=
\begin{pmatrix}
a_{11}&a_{21}&\cdots&a_{m1}\\
a_{12}&a_{22}&\cdots&a_{m2}\\
\vdots&\vdots&\ddots&\vdots\\
a_{1n}&a_{2n}&\cdots&a_{mn}
\end{pmatrix}
\in M_{n\times m}(R),
\qquad
(A^t)_{ij}=a_{ji}.
\]
Thus transpose visibly interchanges rows and columns.

For \(A=(a_{ij})\in M_{m\times n}(F)\) and a column
\[
x=\begin{pmatrix}x_1\\\vdots\\x_n\end{pmatrix},
\]
define matrix--vector multiplication by
\[
(Ax)_i=\sum_{j=1}^n a_{ij}x_j.
\]
For instance,
\[
\begin{pmatrix}2&-1\\3&4\end{pmatrix}
\begin{pmatrix}5\\2\end{pmatrix}
=\begin{pmatrix}8\\23\end{pmatrix}.
\]

For
\[
A=(a_{ij})\in M_{m\times n}(R),
\qquad
B=(b_{jk})\in M_{n\times r}(R),
\]
define
\[
AB=(c_{ik})\in M_{m\times r}(R),
\qquad
c_{ik}=\sum_{j=1}^n a_{ij}b_{jk}.
\]
This is row-by-column multiplication, and its dimension rule is
\[
(m\times n)(n\times r)\longrightarrow(m\times r).
\]
Associativity follows by distributing the finite sums; no commutativity of
\(R\) is required. The \emph{identity matrix}
\[
I_n=(\delta_{ij})\in M_n(R)
\]
has entries \(1\) on the diagonal and \(0\) elsewhere. A square matrix \(A\)
is \emph{invertible} if some square matrix \(B\) satisfies \(AB=BA=I_n\);
then \(B\) is unique and is denoted by \(A^{-1}\).

\begin{proposition}[Transpose identities]
For compatible matrices over a commutative ring,
\[
(A^t)^t=A,
\qquad
(A+B)^t=A^t+B^t,
\qquad
(rA)^t=rA^t,
\qquad
(AB)^t=B^tA^t.
\]
\end{proposition}
\begin{proof}
The first three identities follow by comparing entries. For the product,
\[
\bigl((AB)^t\bigr)_{ij}
=(AB)_{ji}
=\sum_k a_{jk}b_{ki}
=\sum_k b_{ki}a_{jk}
=(B^tA^t)_{ij}.
\]
Commutativity is used in the third equality.
\end{proof}

Now let \(T:V\to W\), let
\[
\beta=(v_1,\ldots,v_n)
\]
be an ordered basis of \(V\), and let
\[
\gamma=(w_1,\ldots,w_m)
\]
be an ordered basis of \(W\). Write
\[
\begin{aligned}
T(v_1)&=a_{11}w_1+a_{21}w_2+\cdots+a_{m1}w_m,\\
T(v_2)&=a_{12}w_1+a_{22}w_2+\cdots+a_{m2}w_m,\\
&\ \vdots\\
T(v_n)&=a_{1n}w_1+a_{2n}w_2+\cdots+a_{mn}w_m.
\end{aligned}
\]
The \emph{matrix of \(T\) from \(\beta\) to \(\gamma\)} is
\[
[T]^\gamma_\beta=
\begin{pmatrix}
a_{11}&a_{12}&\cdots&a_{1n}\\
a_{21}&a_{22}&\cdots&a_{2n}\\
\vdots&\vdots&\ddots&\vdots\\
a_{m1}&a_{m2}&\cdots&a_{mn}
\end{pmatrix}.
\]
Equivalently,
\[
\boxed{\text{the \(j\)-th column of \([T]^\gamma_\beta\) is
\([T(v_j)]_\gamma\).}}
\]
Thus the matrix depends on the chosen ordered bases, though the linear map
does not. If
\[
v=x_1v_1+\cdots+x_nv_n,
\]
then linearity and the displayed expansions give
\[
T(v)
=\sum_{i=1}^m\left(\sum_{j=1}^n a_{ij}x_j\right)w_i.
\]
The inner sum is the \(i\)-th entry of
\([T]^\gamma_\beta[v]_\beta\), and hence
\[
[T(v)]_\gamma=[T]^\gamma_\beta[v]_\beta. \tag{6.1}
\]
For \(S:W\to X\), with an ordered basis \(\delta\) of \(X\), this immediately
explains the order of matrix multiplication:
\[
[S\circ T]^\delta_\beta=[S]^\delta_\gamma[T]^\gamma_\beta.
\]
Formula (6.1) also shows that a square matrix representing an endomorphism is
invertible exactly when that endomorphism is an automorphism.

\begin{example}[Constructing a representing matrix]
Let
\[
T:\mathbb R^2\longrightarrow\mathbb R^2,
\qquad
T(x,y)=(x+y,x-y),
\]
and choose the ordered bases
\[
\beta=((1,1),(1,0)),
\qquad
\gamma=((1,0),(1,1)).
\]
Then
\[
T(1,1)=(2,0)=2(1,0)+0(1,1),
\qquad
T(1,0)=(1,1)=0(1,0)+1(1,1).
\]
The coordinate columns therefore give
\[
[T]^\gamma_\beta=
\begin{pmatrix}2&0\\0&1\end{pmatrix}.
\]
For
\[
v=2(1,1)-(1,0)=(1,2),
\]
we have
\[
[v]_\beta=
\begin{pmatrix}2\\-1\end{pmatrix},
\qquad
[T(v)]_\gamma=
\begin{pmatrix}4\\-1\end{pmatrix}
=
\begin{pmatrix}2&0\\0&1\end{pmatrix}
\begin{pmatrix}2\\-1\end{pmatrix}.
\]
\end{example}

If \(\beta'\) is another ordered basis of \(V\), define the
change-of-coordinate matrix, with its direction displayed, by
\(P=[I_V]^\beta_{\beta'}\). Formula (6.1) for the identity says
\[
[v]_\beta=P[v]_{\beta'};
\]
thus \(P\) converts \(\beta'\)-coordinates into \(\beta\)-coordinates. The
reverse change-of-coordinate matrix
\(Q=[I_V]^{\beta'}_\beta\) satisfies \(PQ=QP=I\), so \(P^{-1}=Q\).
For an endomorphism \(T\), applying the composition formula to the two
identity maps gives
\[
[T]^{\beta'}_{\beta'}=P^{-1}[T]^\beta_\beta P.
\]
Matrices related in this way are \emph{similar}; they represent one operator
in different ordered bases.

Finally, the matrices \(E_{ij}\), having a single \(1\) in position
\((i,j)\), form a basis of \(M_{m\times n}(F)\). Hence this matrix space has
dimension \(mn\). For fixed ordered bases, \(T\mapsto[T]^\gamma_\beta\) is
a linear bijection
\[
\operatorname{Hom}_F(V,W)\longrightarrow M_{m\times n}(F)
\]
by the column construction above; this identification depends on the chosen
bases. In particular,
\(\dim_F\mathcal{L}_F(V)=(\dim_FV)^2\).

\begin{example}[A change of basis]
Let \(T:\mathbb R^2\to\mathbb R^2\) be given by
\(T(x,y)=(3x+y,y)\). In the standard basis its matrix is
\(A=\begin{pmatrix}3&1\\0&1\end{pmatrix}\). Put
\(v_1=(1,0)\), \(v_2=(1,-2)\), and let \(\beta=(v_1,v_2)\). The
change-of-coordinate matrix has these vectors as columns,
\[P=\begin{pmatrix}1&1\\0&-2\end{pmatrix},\qquad
P^{-1}=\begin{pmatrix}1&\tfrac12\\0&-\tfrac12\end{pmatrix}.
\]
Therefore
\[ [T]_\beta=P^{-1}AP=\begin{pmatrix}3&0\\0&1\end{pmatrix}. \]
The diagonal entries were not produced by a new operator: the chosen basis
consists of eigenvectors. This computation is the concrete meaning of
similarity and anticipates diagonalization.
\end{example}

\section{Determinants}\label{sec:determinants}
An endomorphism of an \(n\)-dimensional space either preserves oriented
\(n\)-dimensional volume up to a scalar or collapses some direction. The
determinant is that scalar. Its most useful definition is therefore
structural, not computational. Permutations and their signs, used below, were
proved in \hyperref[chap:groups-foundations]{Chapter~2}.

\begin{definition}
Let \(V\) be an \(n\)-dimensional vector space over \(F\), with ordered basis
\(\beta=(v_1,\ldots,v_n)\). A map \(\omega:V^n\to F\) is
\emph{\(n\)-multilinear} if it is linear in each argument separately, and
\emph{alternating} if it vanishes whenever two arguments agree. The
\emph{determinant form associated with \(\beta\)} is an alternating
\(n\)-multilinear map \(\omega_\beta\) satisfying
\[
\omega_\beta(v_1,\ldots,v_n)=1.
\]
\end{definition}

Multilinearity models how volume changes when one edge is changed; alternation
models the collapse caused by two equal directions; normalization fixes the
unit oriented parallelepiped determined by the chosen basis.

\begin{lemma}[Alternating sign-change lemma]\label{lem:alternating-sign-change}
Let \(\omega:V^n\to F\) be alternating and multilinear. Interchanging any
two arguments changes its value by a sign. More precisely, if every argument
other than those in two fixed positions is held fixed, then
\[
\omega(\ldots,x,\ldots,y,\ldots)
=-\omega(\ldots,y,\ldots,x,\ldots).
\]
\end{lemma}
\begin{proof}
Alternation gives
\[
0=\omega(\ldots,x+y,\ldots,x+y,\ldots).
\]
Expanding the two displayed positions by multilinearity gives four terms.
The terms with two copies of \(x\) or two copies of \(y\) vanish by
alternation. Thus
\[
0=\omega(\ldots,x,\ldots,y,\ldots)
 +\omega(\ldots,y,\ldots,x,\ldots),
\]
which is the asserted identity.
\end{proof}

Consequently, each transposition of arguments contributes a factor of
\(-1\). Chapter~2 proves that the parity of a permutation is well-defined,
so a permutation \(\sigma\in S_n\) changes the value of \(\omega\) by the
factor \(\operatorname{sgn}(\sigma)\). The lemma supplies the effect on an
alternating multilinear map; the earlier permutation theory supplies the
rule by which the signs of successive transpositions combine.

\begin{theorem}[Existence and uniqueness of the determinant form]\label{thm:det-form}
There is exactly one determinant form \(\omega_\beta\).
\end{theorem}
\begin{proof}
We first prove uniqueness. Write
\[
x_j=\sum_{i=1}^n a_{ij}v_i\qquad(1\leq j\leq n).
\]
Repeated multilinearity forces
\[
\omega(x_1,\ldots,x_n)
=\sum_{i_1,\ldots,i_n}
a_{i_1,1}\cdots a_{i_n,n}\,
\omega(v_{i_1},\ldots,v_{i_n}).
\]
Alternation kills every term with repeated indices. The surviving tuples are
exactly
\[
(i_1,\ldots,i_n)=(\sigma(1),\ldots,\sigma(n)),
\qquad \sigma\in S_n.
\]
By Lemma~\ref{lem:alternating-sign-change}, each adjacent transposition of
arguments contributes a factor of \(-1\). The Chapter~2 parity theory then
shows that the term indexed by \(\sigma\) has the factor
\(\operatorname{sgn}(\sigma)\), and normalization gives
\[
\omega_\beta(x_1,\ldots,x_n)
=\sum_{\sigma\in S_n}\operatorname{sgn}(\sigma)
a_{\sigma(1),1}\cdots a_{\sigma(n),n}. \tag{6.2}\label{eq:leibniz-form}
\]
Thus there can be at most one such map.

Conversely, define the right side of \eqref{eq:leibniz-form} to be
\(\omega_\beta(x_1,\ldots,x_n)\). Each coefficient in one column occurs
linearly, so the expression is multilinear in the corresponding argument.
If \(x_r=x_s\), pair a summand indexed by \(\sigma\) with the summand obtained
by exchanging \(\sigma(r)\) and \(\sigma(s)\). The products agree because
columns \(r\) and \(s\) agree, while the signs are opposite; every summand
therefore cancels. Finally, for \((x_1,\ldots,x_n)=(v_1,\ldots,v_n)\), only
the identity permutation contributes, and the value is \(1\). Hence the
candidate is a determinant form.
\end{proof}

Formula \eqref{eq:leibniz-form}, the \emph{Leibniz formula}, is consequently
not a separate definition: it is the forced coordinate expression for the
unique normalized alternating multilinear form.

\begin{definition}
For \(A\in M_n(F)\), with columns \(A_1,\ldots,A_n\), define
\[
\det(A)=\omega_e(A_1,\ldots,A_n),
\]
where \(e\) is the standard ordered basis of \(F^n\). Thus
\[
\det(A)=
\sum_{\sigma\in S_n}\operatorname{sgn}(\sigma)
a_{\sigma(1),1}\cdots a_{\sigma(n),n}.
\]
This is a column convention. The row convention follows below from transpose
invariance.
\end{definition}

\begin{definition}
For \(T\in\mathcal{L}_F(V)\), define \(\det(T)\) by
\[
\omega_\beta(Tx_1,\ldots,Tx_n)
=\det(T)\,\omega_\beta(x_1,\ldots,x_n)
\qquad(x_1,\ldots,x_n\in V).
\]
\end{definition}

The left side is alternating and multilinear in the \(x_i\). The expansion
used in the proof of Theorem~\ref{thm:det-form} shows that every such form is
its value on \(\beta\) times \(\omega_\beta\), so the displayed scalar exists
and is unique; evaluating at \(\beta\) shows that it equals
\(\omega_\beta(Tv_1,\ldots,Tv_n)\). The identity itself holds for every
tuple, so evaluating in another ordered basis proves that \(\det(T)\) is
independent of \(\beta\). Expanding the columns \(Tv_j\) in \(\beta\) gives
\[
\det(T)=\det([T]^\beta_\beta). \tag{6.3}\label{eq:operator-matrix-det}
\]
This separates the intrinsic operator from its coordinate matrix.

\begin{theorem}[Basic determinant identities]\label{thm:det-identities}
For endomorphisms \(S,T\) of a finite-dimensional vector space over a field,
\[
\det(I)=1,\qquad
\det(S\circ T)=\det(S)\det(T).
\]
For square matrices over \(F\),
\[
\det(AB)=\det(A)\det(B),
\qquad
\det(A^t)=\det(A).
\]
\end{theorem}
\begin{proof}
The defining identity for determinants gives
\[
\omega_\beta(STx_1,\ldots,STx_n)
=\det(S)\det(T)\omega_\beta(x_1,\ldots,x_n),
\]
which proves multiplicativity and \(\det(I)=1\). Matrix multiplicativity
follows from \eqref{eq:operator-matrix-det} and the composition formula
proved in \S6.11. Replacing \(\sigma\) by \(\sigma^{-1}\) in
\eqref{eq:leibniz-form} gives \(\det(A^t)=\det(A)\).
\end{proof}

Thus adding a multiple of one column to another leaves the determinant
unchanged, scaling a column scales the determinant by the same scalar, and
interchanging two columns changes its sign; these are immediate consequences
of multilinearity and alternation. Transpose invariance supplies the
corresponding row operations.

\begin{definition}
For \(A=(a_{ij})\in M_n(F)\), let \(A_{ij}\) be the matrix obtained by
deleting row \(i\) and column \(j\). Its \emph{cofactor} is
\[
C_{ij}=(-1)^{i+j}\det(A_{ij}).
\]
\end{definition}

\begin{theorem}[Laplace expansion and adjugate]\label{thm:adjugate}
For every row \(i\) and column \(j\),
\[
\det A=\sum_{k=1}^n a_{ik}C_{ik},
\qquad
\det A=\sum_{k=1}^n a_{kj}C_{kj}.
\]
If \(\operatorname{adj}(A)=(C_{ji})\), then
\[
A\operatorname{adj}(A)=\operatorname{adj}(A)A=(\det A)I_n. \tag{6.4}\label{eq:adjugate}
\]
\end{theorem}
\begin{proof}
Group the terms of the Leibniz formula according to the column selected in
row \(i\). Moving that selected entry to the first position requires
\(i+j\) parity changes, giving the first expansion; transpose invariance
gives the column version. The \((i,j)\)-entry of
\(A\operatorname{adj}(A)\) is the row-\(i\) cofactor expansion if \(i=j\).
If \(i\ne j\), it is the same expansion after replacing row \(j\) by row
\(i\), which is zero by alternation. The column calculation proves the
identity on the other side.
\end{proof}

\begin{corollary}[Invertibility criterion over a field]\label{cor:field-det-invertibility}
For \(A\in M_n(F)\), equivalently for an endomorphism of \(V\),
\[
A\text{ is invertible}\quad\Longleftrightarrow\quad\det(A)\ne0.
\]
\end{corollary}
\begin{proof}
If \(A\) is invertible, determinant multiplicativity gives
\[
1=\det(A)\det(A^{-1}).
\]
Conversely, if \(\det(A)\ne0\), then
\(\det(A)^{-1}\operatorname{adj}(A)\) is an inverse by
\eqref{eq:adjugate}.
\end{proof}

\begin{example}
For a \(2\times2\) matrix,
\[
\det\begin{pmatrix}a&b\\c&d\end{pmatrix}=ad-bc.
\]
Thus \(\det\begin{pmatrix}2&1\\4&2\end{pmatrix}=0\), while
\[
\det\begin{pmatrix}2&1\\1&1\end{pmatrix}=1,
\qquad
\begin{pmatrix}2&1\\1&1\end{pmatrix}^{-1}
=\begin{pmatrix}1&-1\\-1&2\end{pmatrix}.
\]
Laplace expansion and the adjugate are computational consequences of the
same structural determinant form, not competing definitions.
\end{example}

\subsection{Determinants over commutative rings}
The Leibniz expression uses only addition, multiplication, commutativity,
and the elements \(1\) and \(-1\). It therefore has a useful scope beyond
fields.

\begin{theorem}[Determinants over commutative rings]\label{thm:ring-determinants}
Let \(R\) be a commutative ring with identity. For \(A\in M_n(R)\), define
\[
\det(A)=
\sum_{\sigma\in S_n}\operatorname{sgn}(\sigma)
a_{\sigma(1),1}\cdots a_{\sigma(n),n}.
\]
On \(R^n\), this is the unique normalized alternating \(n\)-multilinear
form with respect to the standard ordered basis. The identities
\[
\det(I)=1,\qquad
\det(AB)=\det(A)\det(B),\qquad
\det(A^t)=\det(A)
\]
and the Laplace and adjugate identities \eqref{eq:adjugate} remain valid.
\end{theorem}
\begin{proof}
The uniqueness, cancellation, and multilinearity arguments in
Theorem~\ref{thm:det-form} use no division, only commutativity and the sign
of a permutation. The structural proof of multiplicativity and the cofactor
proof use the same operations, and hence carry over word for word.
\end{proof}

If \(P\) is an invertible change-of-basis matrix of a finite free
\(R\)-module \(F\), then
\[
\det(P)\det(P^{-1})=1.
\]
Consequently, for an endomorphism \(T:F\to F\), the value
\[
\det(T)=\det([T]^\beta_\beta)
\]
is independent of the ordered basis \(\beta\): two matrices are related by
\(P^{-1}[T]^\beta_\beta P\), and determinant multiplicativity cancels the
two unit determinants. Thus the preceding field construction extends exactly
to finite free modules, not to arbitrary finitely generated modules.

\begin{corollary}[Invertibility over a commutative ring]
For \(A\in M_n(R)\),
\[
A\text{ is invertible}\quad\Longleftrightarrow\quad\det(A)\in R^\times.
\]
\end{corollary}
\begin{proof}
An inverse of \(A\) makes \(\det(A)\det(A^{-1})=1\). Conversely, if
\(\det(A)\) is a unit, \((\det A)^{-1}\operatorname{adj}(A)\) is an inverse
by \eqref{eq:adjugate}.
\end{proof}

Over a field, nonzero and unit are equivalent. Over \(\mathbb Z\), they are
not: the \(1\times1\) matrix \((2)\) has nonzero determinant but is not
invertible. In particular, the field criterion involving linear dependence
must not be transferred unqualified to general commutative rings. For
\(A\in M_n(R)\), the polynomial
\(\chi_A(t)=\det(tI-A)\in R[t]\) is nevertheless defined, and the adjugate
identity over \(R[t]\) gives \(\chi_A(A)=0\). Chapter~7 uses this argument in
its operator-theoretic field setting.

\section{Eigenvalues, Eigenvectors, and Invariant Subspaces}
An eigenvector is a direction whose direction is not changed by an operator:
the operator merely scales it. Such directions are the first indication that
a complicated transformation might have a simple coordinate system.

\begin{definition}
For \(T:V\to V\), a nonzero \(v\) is an \emph{eigenvector} with
\emph{eigenvalue} \(\lambda\in F\) if \(Tv=\lambda v\). The
\emph{eigenspace} is \(E_\lambda=\ker(T-\lambda I)\). A subspace
\(W\leq V\) is \emph{\(T\)-invariant} if \(T(W)\subseteq W\).
\end{definition}

Because eigenspaces are kernels, they are subspaces before any dimension is
attached to them. Now let \(\beta\) be an ordered basis and
\(A=[T]^\beta_\beta\). Define the \emph{characteristic polynomial} by
\[\chi_T(x)=\det(xI-A)\in F[x].\]
This definition is independent of \(\beta\). Indeed, if
\(A'=P^{-1}AP\), then, in the polynomial ring \(F[x]\),
\[
xI-A'=P^{-1}(xI-A)P.
\]
Determinant multiplicativity, already proved above, gives
\(\det(xI-A')=\det(P^{-1})\det(xI-A)\det(P)=\det(xI-A)\).

\begin{proposition}
For \(\lambda\in F\), \(\lambda\) is an eigenvalue of \(T\) if and only if
\(\chi_T(\lambda)=0\).
\end{proposition}
\begin{proof}
The equation \(Tv=\lambda v\) for a nonzero \(v\) says that
\(T-\lambda I\) is not injective. In finite dimension it is therefore not
invertible. Its matrix is \(A-\lambda I\), so the determinant criterion
above says precisely that \(\det(A-\lambda I)=0\), equivalently
\(\det(\lambda I-A)=0\). This is \(\chi_T(\lambda)=0\). The reverse chain
of equivalences proves the converse.
\end{proof}

The eigenspace is a subspace because it is a kernel. It can have dimension
larger than one: for the identity operator, the unique eigenspace is all of
\(V\). By contrast, a rotation of the real plane through a nonzero angle less
than \(\pi\) has no real eigenvector at all. The field therefore matters:
the same rotation has complex eigenvalues, a point returned to in Chapter~7.

\section{Diagonalization}
The best possible matrix of an operator is diagonal: each coordinate evolves
independently. The question is therefore not whether a matrix happens to
look diagonal, but whether the underlying operator has a basis made of
eigenvectors.

\begin{theorem}[Diagonalization criterion]
For an endomorphism \(T:V\to V\), the following are equivalent:
\begin{enumerate}
\item \(T\) is represented by a diagonal matrix in some basis;
\item \(V\) has a basis of eigenvectors of \(T\);
\item
\[
V=\bigoplus_\lambda E_\lambda,
\]
where the sum ranges over the
eigenvalues that occur.
\end{enumerate}
\end{theorem}
\begin{proof}
The columns of a diagonal matrix say exactly that its basis vectors are
eigenvectors, proving (1) iff (2). Eigenvectors belonging to distinct
eigenvalues are independent. For if a shortest nontrivial relation
\[
\sum_{i=1}^rc_iv_i=0
\]
involved distinct eigenvalues \(\lambda_i\), then
applying \(T-\lambda_rI\) gives a shorter relation
\[
\sum_{i<r}c_i(\lambda_i-\lambda_r)v_i=0;
\]
its coefficients are nonzero
whenever the original ones are, a contradiction. Thus the sum of the
eigenspaces is direct. A basis of eigenvectors spans precisely this direct
sum, proving (2) iff (3).
\end{proof}

In particular, if \(\dim V=n\) and \(T\) has \(n\) distinct eigenvalues,
then it is diagonalizable. Diagonalization is powerful but not universal. The
next chapter explains the additional structure that remains when no eigenbasis
exists.

\begin{example}[Three first tests for diagonalization]
On \(\mathbb R^2\), the operator with matrix
\(\begin{pmatrix}3&1\\0&1\end{pmatrix}\) has distinct eigenvalues \(3\) and
\(1\), so the preceding worked calculation diagonalizes it. The identity
operator has the repeated eigenvalue \(1\), yet is diagonalizable because
every nonzero vector is an eigenvector. In contrast,
\[J=\begin{pmatrix}1&1\\0&1\end{pmatrix}\]
has characteristic polynomial \((x-1)^2\), but
\(\ker(J-I)\) is one-dimensional. It has only one independent eigenvector
and therefore is not diagonalizable. Thus a split characteristic polynomial
does not by itself supply an eigenbasis; the missing information is precisely
what the module structure in Chapter~7 records.
\end{example}

\section{Dual Spaces and Dual Maps}\label{sec:dual-spaces}
Coordinates record a vector after a basis has been chosen.  There is another
useful way to probe a vector: apply every scalar-valued linear measurement to
it.  This viewpoint is intrinsic, and it is the starting point for
bilinear forms and alternating forms.

\begin{definition}
A \emph{linear functional} on an \(F\)-vector space \(V\) is a linear map
\(\varphi:V\to F\).  The \emph{algebraic dual space} of \(V\) is
\[
V^\vee=\operatorname{Hom}_F(V,F).
\]
\end{definition}

The notation \(V^\vee\) is deliberately different from the notation
\(T^*\) used for the inner-product adjoint in Chapter~7. The construction
below is algebraic and requires no inner product.

\begin{proposition}[Dual basis]\label{prop:dual-basis}
Let \(\beta=(v_1,\ldots,v_n)\) be an ordered basis of a finite-dimensional
vector space \(V\).  There are unique linear functionals
\[
v_i^\vee:V\longrightarrow F
\]
such that
\[
v_i^\vee(v_j)=\delta_{ij}.
\]
Moreover,
\[
\beta^\vee=(v_1^\vee,\ldots,v_n^\vee)
\]
is a basis of \(V^\vee\).  In particular,
\[
\dim V^\vee=\dim V.
\]
\end{proposition}
\begin{proof}
For
\[
v=\sum_{j=1}^na_jv_j,
\]
define \(v_i^\vee(v)=a_i\).  The uniqueness of basis coordinates proves
that this is well-defined and linear, and gives the displayed values.
If
\[
\sum_{i=1}^nc_iv_i^\vee=0,
\]
evaluation at \(v_j\) gives \(c_j=0\), so the functionals are independent.
For \(\varphi\in V^\vee\), linearity gives
\[
\varphi=\sum_{i=1}^n\varphi(v_i)v_i^\vee.
\]
Thus they span and form a basis.
\end{proof}

Consequently \(V\) and \(V^\vee\) are isomorphic when \(V\) is
finite-dimensional, but this proof used the chosen basis.  In general there
is no preferred isomorphism
\[
V\cong V^\vee.
\]
An inner product can supply one, but that is additional structure.

\begin{definition}
For a linear map
\[
T:V\longrightarrow W,
\]
its \emph{dual map}, also called its algebraic transpose, is
\[
T^\vee:W^\vee\longrightarrow V^\vee,
\qquad
T^\vee(\varphi)=\varphi\circ T.
\]
\end{definition}

The direction reverses because a measurement on \(W\) can be pulled back
along \(T\) to a measurement on \(V\).

\begin{proposition}[Functorial identities and transpose matrix]
For linear maps \(T:V\to W\) and \(S:W\to U\),
\[
(\operatorname{id}_V)^\vee=\operatorname{id}_{V^\vee},
\qquad
(S\circ T)^\vee=T^\vee\circ S^\vee.
\]
If \(\beta\) and \(\gamma\) are ordered bases of \(V\) and \(W\), then
\[
[T^\vee]^{\beta^\vee}_{\gamma^\vee}
=\left([T]^\gamma_\beta\right)^t.
\]
\end{proposition}
\begin{proof}
For \(\varphi\in V^\vee\),
\[
(\operatorname{id}_V)^\vee(\varphi)
=\varphi\circ\operatorname{id}_V=\varphi.
\]
Likewise, for \(\psi\in U^\vee\),
\[
(S\circ T)^\vee(\psi)
=\psi\circ S\circ T
=T^\vee(S^\vee(\psi)).
\]
Write
\[
T(v_j)=\sum_{i=1}^ma_{ij}w_i,
\qquad
[T]^\gamma_\beta=(a_{ij}).
\]
For the \(i\)-th dual basis vector of \(W^\vee\),
\[
T^\vee(w_i^\vee)(v_j)=w_i^\vee(T(v_j))=a_{ij}.
\]
Hence
\[
T^\vee(w_i^\vee)=\sum_{j=1}^na_{ij}v_j^\vee.
\]
Its matrix therefore has \((j,i)\)-entry \(a_{ij}\), which is the transpose
of the matrix of \(T\).
\end{proof}

\begin{definition}
The \emph{evaluation map} is
\[
\iota_V:V\longrightarrow V^{\vee\vee},
\qquad
\iota_V(v)(\varphi)=\varphi(v).
\]
\end{definition}

\begin{proposition}[Double dual]\label{prop:double-dual}
The evaluation map \(\iota_V\) is linear and injective.  If \(V\) is
finite-dimensional, it is an isomorphism.
\end{proposition}
\begin{proof}
For \(a,b\in F\), \(v,w\in V\), and \(\varphi\in V^\vee\),
\[
\iota_V(av+bw)(\varphi)
=\varphi(av+bw)
=a\iota_V(v)(\varphi)+b\iota_V(w)(\varphi),
\]
so \(\iota_V\) is linear.  If \(v\ne0\), extend \(v\) to a basis and take
the corresponding first coordinate functional.  It does not vanish on
\(v\), so \(\iota_V(v)\ne0\).  Thus \(\iota_V\) is injective.  In finite
dimension, Proposition~\ref{prop:dual-basis} gives
\[
\dim V^{\vee\vee}=\dim V,
\]
so an injective map between these finite-dimensional spaces is surjective.
\end{proof}

Thus finite-dimensional vector spaces have a canonical isomorphism
\[
V\cong V^{\vee\vee},
\]
even though an isomorphism with the first dual is usually noncanonical.
This canonical character is expressed by the naturality identity
\begin{equation}\label{eq:double-dual-naturality}
\iota_W\circ T=T^{\vee\vee}\circ\iota_V
\qquad(T:V\to W).
\end{equation}
Indeed, for \(v\in V\) and \(\varphi\in W^\vee\),
\[
\begin{aligned}
\bigl(T^{\vee\vee}\iota_V(v)\bigr)(\varphi)
&=\iota_V(v)(\varphi\circ T)\\
&=(\varphi\circ T)(v)\\
&=\varphi(Tv)\\
&=\iota_W(Tv)(\varphi).
\end{aligned}
\]
Thus evaluation commutes with every linear map and requires no basis choice.

\begin{definition}
For a subspace \(W\leq V\), its \emph{annihilator} is
\[
W^\circ=\{\varphi\in V^\vee:\varphi(w)=0\text{ for every }w\in W\}.
\]
\end{definition}

\begin{proposition}
If \(V\) is finite-dimensional, then
\[
\dim W^\circ=\dim V-\dim W.
\]
\end{proposition}
\begin{proof}
Choose a basis \(w_1,\ldots,w_r\) of \(W\) and extend it to a basis
\[
(w_1,\ldots,w_r,u_1,\ldots,u_s)
\]
of \(V\).  Its dual basis shows that \(W^\circ\) has basis
\[
(u_1^\vee,\ldots,u_s^\vee).
\]
Since \(r+s=\dim V\), the result follows.
\end{proof}

\begin{theorem}[Dual maps, kernels, and images]\label{thm:dual-map-kernel-image}
Let \(T:V\to W\) be a linear map between finite-dimensional vector spaces.
Then
\[
\ker(T^\vee)=(\operatorname{im}T)^\circ,
\qquad
\operatorname{im}(T^\vee)=(\ker T)^\circ.
\]
Consequently,
\[
\operatorname{rank}(T^\vee)=\operatorname{rank}(T).
\]
\end{theorem}
\begin{proof}
For \(\varphi\in W^\vee\), every equivalence is direct:
\[
\begin{aligned}
\varphi\in\ker(T^\vee)
&\Longleftrightarrow T^\vee\varphi=0\\
&\Longleftrightarrow \varphi\circ T=0\\
&\Longleftrightarrow \varphi\text{ vanishes on }\operatorname{im}T\\
&\Longleftrightarrow \varphi\in(\operatorname{im}T)^\circ.
\end{aligned}
\]
Hence, using the annihilator dimension formula and rank--nullity,
\[
\begin{aligned}
\operatorname{rank}(T^\vee)
&=\dim W^\vee-\dim\ker(T^\vee)\\
&=\dim W-\bigl(\dim W-\dim\operatorname{im}T\bigr)\\
&=\operatorname{rank}T.
\end{aligned}
\]
If \(\varphi=T^\vee\psi=\psi\circ T\) and \(v\in\ker T\), then
\(\varphi(v)=\psi(Tv)=0\).  Thus
\(\operatorname{im}(T^\vee)\subseteq(\ker T)^\circ\).  The two spaces have
the same dimension, because
\[
\dim\operatorname{im}(T^\vee)=\operatorname{rank}T
=\dim V-\dim\ker T=\dim(\ker T)^\circ.
\]
The containment is therefore equality.
\end{proof}

In dual bases,
\([T^\vee]=[T]^t\).  The theorem therefore yields
\[
\operatorname{rank}(A^t)=\operatorname{rank}(A),
\]
giving the conceptual dual-space explanation of the row-rank equals
column-rank theorem proved earlier by elementary matrix methods.

\begin{proposition}[Quotient dual]\label{prop:quotient-dual}
If \(W\leq V\) and \(\pi:V\to V/W\) is the quotient map, then
\[
\pi^\vee:(V/W)^\vee\longrightarrow V^\vee
\]
is injective with image \(W^\circ\).  Hence
\[
(V/W)^\vee\cong W^\circ.
\]
\end{proposition}
\begin{proof}
Because \(\pi\) is surjective, \(\psi\circ\pi=0\) implies \(\psi=0\), so
\(\pi^\vee\) is injective.  Theorem~\ref{thm:dual-map-kernel-image} gives
\[
\operatorname{im}(\pi^\vee)=(\ker\pi)^\circ=W^\circ.
\]
\end{proof}

\begin{proposition}[Annihilator identities]\label{prop:annihilator-identities}
For subspaces \(U,W\leq V\), with \(V\) finite-dimensional,
\[
U\subseteq W\Longrightarrow W^\circ\subseteq U^\circ,
\qquad
(U+W)^\circ=U^\circ\cap W^\circ,
\]
\[
(U\cap W)^\circ=U^\circ+W^\circ,
\qquad
W^{\circ\circ}=W
\]
under the evaluation identification \(V\cong V^{\vee\vee}\).
\end{proposition}
\begin{proof}
A functional vanishing on \(W\) certainly vanishes on \(U\subseteq W\),
which proves reversal of inclusions.  A functional vanishes on \(U+W\)
exactly when it vanishes on both summands, proving the second identity.

Every element of \(U^\circ+W^\circ\) vanishes on \(U\cap W\), so
\[
U^\circ+W^\circ\subseteq(U\cap W)^\circ.
\]
The dimension formula for sums and intersections gives
\[
\begin{aligned}
\dim(U^\circ+W^\circ)
&=\dim U^\circ+\dim W^\circ-\dim(U^\circ\cap W^\circ)\\
&=(\dim V-\dim U)+(\dim V-\dim W)-\dim(U+W)^\circ\\
&=\dim V-\dim(U\cap W),
\end{aligned}
\]
which is the dimension of \((U\cap W)^\circ\); hence equality holds.
Finally, evaluation sends every \(w\in W\) to a functional on \(V^\vee\)
that vanishes on \(W^\circ\), so \(W\subseteq W^{\circ\circ}\).  Both spaces
have dimension \(\dim W\), and therefore they are equal.
\end{proof}

\section{Bilinear and Quadratic Forms}\label{sec:bilinear-quadratic-forms}
Duality turns a vector into a linear measurement.  A bilinear form instead
assigns a scalar to an ordered pair of vectors, linearly in each argument.
This is the algebraic framework behind dot products, but it does not include
positivity or conjugation.

\begin{definition}
A \emph{bilinear form} on an \(F\)-vector space \(V\) is a map
\[
B:V\times V\longrightarrow F
\]
that is linear in each variable.  It is \emph{symmetric} if
\[
B(x,y)=B(y,x),
\]
\emph{skew-symmetric} if
\[
B(x,y)=-B(y,x),
\]
and \emph{alternating} if
\[
B(v,v)=0
\qquad(v\in V).
\]
\end{definition}

Every alternating bilinear form is skew-symmetric, by expanding
\[
0=B(x+y,x+y).
\]
Conversely, a skew-symmetric form is alternating only when the
characteristic is not \(2\), in the sense of
Definition~\ref{def:ring-characteristic}. This distinction is essential in
characteristic \(2\).
In every ordered basis \(\beta\),
\[
B\text{ symmetric}\Longleftrightarrow[B]_\beta^t=[B]_\beta,
\qquad
B\text{ skew-symmetric}\Longleftrightarrow[B]_\beta^t=-[B]_\beta.
\]
Alternating always implies the second condition.  If
\(\operatorname{char}F\ne2\), the second condition forces
\(2B(v,v)=0\) and hence \(B(v,v)=0\).  In characteristic \(2\), however,
\(-A=A\), so the transpose equation merely says that \(A\) is symmetric; an
alternating form must additionally have zero diagonal.

\begin{proposition}[Matrix and congruence]\label{prop:bilinear-matrix}
For an ordered basis \(\beta=(v_1,\ldots,v_n)\), put
\[
[B]_\beta=(B(v_i,v_j)).
\]
Then
\[
B(x,y)=[x]_\beta^t[B]_\beta[y]_\beta.
\]
If
\[
P=[I_V]^\beta_{\beta'},
\]
so that
\[
[x]_\beta=P[x]_{\beta'},
\]
then
\[
[B]_{\beta'}=P^t[B]_\beta P.
\]
\end{proposition}
\begin{proof}
Write
\[
x=\sum_i a_iv_i,
\qquad
y=\sum_j b_jv_j.
\]
By bilinearity,
\[
B(x,y)=\sum_{i,j}a_iB(v_i,v_j)b_j,
\]
which is the stated matrix product.  The columns of \(P\) are the
\(\beta\)-coordinates of the vectors of \(\beta'\).  Substitution of
\[
[x]_\beta=P[x]_{\beta'},
\qquad
[y]_\beta=P[y]_{\beta'}
\]
into the coordinate formula gives
\[
B(x,y)
=[x]_{\beta'}^tP^t[B]_\beta P[y]_{\beta'},
\]
and uniqueness of a matrix representing the form in \(\beta'\) proves the
claim.
\end{proof}

This is \emph{congruence}, not similarity.  For an operator, one input is
converted into output coordinates and change of basis gives
\[
[T]_{\beta'}=P^{-1}[T]_\beta P.
\]
A bilinear form has two vector inputs, so both coordinates are substituted
and the formula has \(P^t\) on the left instead.

\begin{definition}
The \emph{left radical} and \emph{right radical} of \(B\) are, respectively,
\[
\operatorname{rad}_L(B)
=\{v\in V:B(v,w)=0\text{ for every }w\in V\},
\]
\[
\operatorname{rad}_R(B)
=\{w\in V:B(v,w)=0\text{ for every }v\in V\}.
\]
For a general bilinear form these are conceptually distinct.  The associated
linear maps are
\[
\Phi_B:V\longrightarrow V^\vee,
\qquad
\Phi_B(v)=B(v,-).
\]
and
\[
\Psi_B:V\longrightarrow V^\vee,
\qquad
\Psi_B(w)=B(-,w).
\]
\end{definition}

\begin{definition}
The \emph{rank of a bilinear form} is
\[
\operatorname{rank}B:=\operatorname{rank}\Phi_B.
\]
\end{definition}

\begin{proposition}[Rank and radicals of a bilinear form]
For every ordered basis \(\beta\),
\[
\operatorname{rank}B=\operatorname{rank}[B]_\beta.
\]
This number is independent of the basis.  Moreover,
\[
\dim\operatorname{rad}_L(B)=\dim V-\operatorname{rank}B,
\qquad
\dim\operatorname{rad}_R(B)=\dim V-\operatorname{rank}B.
\]
\end{proposition}
\begin{proof}
In \(\beta\) and its dual basis, the matrix of \(\Phi_B\) is
\([B]_\beta^t\).  Equality of row and column rank therefore gives
\[
\operatorname{rank}\Phi_B
=\operatorname{rank}[B]_\beta^t
=\operatorname{rank}[B]_\beta.
\]
Alternatively, basis independence follows directly from
\([B]_{\beta'}=P^t[B]_\beta P\) and invariance of rank under multiplication
by invertible matrices.  Rank--nullity applied to \(\Phi_B\) proves the
left-radical formula.  The right-radical formula follows in the same way from
\(\Psi_B\), whose matrix is \([B]_\beta\).
\end{proof}

\begin{proposition}[Nondegeneracy criterion]
For a finite-dimensional vector space, the following are equivalent:
\begin{enumerate}
\item \(\operatorname{rad}_L(B)=\{0\}\);
\item \(\operatorname{rad}_R(B)=\{0\}\);
\item \([B]_\beta\) is invertible for one, equivalently every, basis
\(\beta\).
\end{enumerate}
When these conditions hold, \(B\) is called \emph{nondegenerate}.
\end{proposition}
\begin{proof}
By their definitions,
\[
\ker\Phi_B=\operatorname{rad}_L(B),
\qquad
\ker\Psi_B=\operatorname{rad}_R(B).
\]
In the basis \(\beta\) and its dual, the matrices of these maps are
\[
[\Phi_B]_{\beta}^{\beta^\vee}=[B]_\beta^t,
\qquad
[\Psi_B]_{\beta}^{\beta^\vee}=[B]_\beta.
\]
Indeed, the \(j\)-th coordinate of \(\Phi_B(v_i)\) is \(B(v_i,v_j)\),
whereas that of \(\Psi_B(v_i)\) is \(B(v_j,v_i)\).  Thus either map is
injective exactly when \([B]_\beta\) is invertible.  Since
\[
\dim V=\dim V^\vee,
\]
injectivity is equivalent to bijectivity for both maps.  Finally, the
congruence formula in Proposition~\ref{prop:bilinear-matrix} shows that
matrix invertibility is independent of the chosen basis.
\end{proof}

Equivalently, in finite dimension,
\[
\begin{aligned}
B\text{ nondegenerate}
&\Longleftrightarrow \operatorname{rank}B=\dim V\\
&\Longleftrightarrow [B]_\beta\text{ is invertible}
\Longleftrightarrow \det[B]_\beta\ne0.
\end{aligned}
\]
A chosen nondegenerate bilinear form therefore supplies the isomorphism
\[
\Phi_B:V\xrightarrow{\ \sim\ }V^\vee,\qquad v\longmapsto B(v,-).
\]
This is precisely the extra structure missing from the bare vector space:
there is no preferred isomorphism \(V\cong V^\vee\), but a nondegenerate form
provides one.

For a symmetric bilinear form, \(B(v,w)=B(w,v)\), so the two radicals
coincide.  In that case we write their common value as \(\operatorname{rad}(B)\).

\begin{definition}
For a symmetric bilinear form \(B\) and a subspace \(W\leq V\), define
\[
W^{\perp_B}=\{v\in V:B(v,w)=0\text{ for every }w\in W\}.
\]
\end{definition}

\begin{proposition}[Orthogonal complements for a bilinear form]
If \(B\) is nondegenerate on finite-dimensional \(V\), then
\[
\dim W^{\perp_B}=\dim V-\dim W.
\]
Moreover,
\[
V=W\oplus W^{\perp_B}
\quad\Longleftrightarrow\quad
B|_W\text{ is nondegenerate}.
\]
\end{proposition}
\begin{proof}
The isomorphism \(\Phi_B:V\to V^\vee\) carries \(W^{\perp_B}\) onto the
annihilator \(W^\circ\).  The annihilator dimension formula gives the first
claim.  Also,
\[
W\cap W^{\perp_B}=\operatorname{rad}(B|_W).
\]
Thus the sum is direct exactly when \(B|_W\) is nondegenerate; in that case
the dimension formula shows that the direct sum has dimension \(\dim V\) and
therefore equals \(V\).  The converse is immediate from the displayed
intersection.
\end{proof}

Unlike the positive-definite case, nondegeneracy of \(B\) on \(V\) alone does
not guarantee this direct sum.  For example, on \(F^2\) with
\[
[B]=\begin{pmatrix}0&1\\1&0\end{pmatrix},
\]
the form is nondegenerate, but \(e_1\) is isotropic and
\(Fe_1=(Fe_1)^{\perp_B}\).  Thus isotropic vectors can make
\(W\cap W^{\perp_B}\ne0\).

\begin{definition}
Assume \(\operatorname{char}F\ne2\), in the sense of
Definition~\ref{def:ring-characteristic}. A quadratic form associated with a
symmetric bilinear form \(B\) is
\[
q(v)=B(v,v).
\]
\end{definition}

The form can be recovered from its quadratic function by polarization:
\[
B(x,y)
=\frac12\bigl(q(x+y)-q(x)-q(y)\bigr).
\]
Characteristic \(2\) requires a separate theory: a quadratic form must not
there be identified blindly with a symmetric bilinear form.

In coordinates, if \(A=[B]_\beta=A^t\) and
\(x=(x_1,\ldots,x_n)^t\), then
\[
q(x)=x^tAx
=\sum_i a_{ii}x_i^2+2\sum_{i<j}a_{ij}x_ix_j.
\]
Conversely, the polarization formula recovers the unique symmetric bilinear
form from \(q\).  If \(x=Py\) and \(P^tAP=D\), then
\[
q(Py)=(Py)^tA(Py)=y^tP^tAPy=y^tDy
=d_1y_1^2+\cdots+d_ny_n^2.
\]
Thus congruence is exactly the coordinate change that diagonalizes a
quadratic form.

\begin{example}[Completing the square]
Over a field of characteristic not \(2\), consider
\[
q(x,y)=x^2+4xy+3y^2.
\]
The cross term disappears visibly:
\[
q(x,y)=(x+2y)^2-y^2.
\]
Set \(u=x+2y\) and \(v=y\), equivalently
\[
\begin{pmatrix}x\\y\end{pmatrix}
=\begin{pmatrix}1&-2\\0&1\end{pmatrix}
\begin{pmatrix}u\\v\end{pmatrix}.
\]
Then \(q=u^2-v^2\).  This is the coordinate mechanism behind diagonalization,
not a classification of all quadratic forms.
\end{example}

\begin{theorem}[Diagonalization of symmetric bilinear forms]
Suppose that \(\operatorname{char}F\ne2\).  Every symmetric bilinear form on
a finite-dimensional \(F\)-vector space has a basis in which its matrix is
diagonal.  Equivalently, its quadratic form has coordinates
\[
q(x_1,\ldots,x_n)=a_1x_1^2+\cdots+a_nx_n^2.
\]
\end{theorem}
\begin{proof}
Induct on \(\dim V\).  If \(B=0\), there is nothing to prove.  Otherwise
some \(v\) satisfies \(B(v,v)\ne0\): if every such value were zero, then
polarization would give \(B=0\).  Let
\[
W=\{w\in V:B(v,w)=0\}.
\]
For \(x\in V\),
\[
x=\frac{B(x,v)}{B(v,v)}v+
\left(x-\frac{B(x,v)}{B(v,v)}v\right),
\]
and the parenthesized vector lies in \(W\).  Thus
\[
V=Fv\oplus W.
\]
The restriction of \(B\) to \(W\) is symmetric, so induction supplies a
diagonalizing basis of \(W\); adjoining \(v\) gives the desired basis.
\end{proof}

Over \(\mathbb R\), inner products are precisely positive-definite symmetric
bilinear forms. Over \(\mathbb C\), the inner products of Chapter~7 are
Hermitian: they are linear in the first variable and conjugate-linear in the
second. Thus Chapter~7 is not merely the complex special case of this
bilinear discussion.


\section*{Exercises}
\begin{exercise}[Basic] Verify directly that abelian groups are precisely
\(\mathbb Z\)-modules, and determine the submodules of
\(\mathbb Z/12\mathbb Z\).\end{exercise}
\begin{exercise}[Core] Prove the universal property of the direct sum for an
arbitrary family of modules.\end{exercise}
\begin{exercise}[Core] Let \(T:F^4\to F^3\) have rank \(2\). Find its
nullity and decide whether it can be injective or surjective.\end{exercise}
\begin{exercise}[Core] Compute the matrix of differentiation on polynomials
of degree at most \(3\), using the basis \(1,x,x^2,x^3\).\end{exercise}
\begin{exercise}[Core] Show that similar matrices have the same characteristic
polynomial and determinant.\end{exercise}
\begin{exercise}[Further] Decide which of the matrices
\(\begin{pmatrix}1&1\\0&1\end{pmatrix}\) and
\(\begin{pmatrix}1&0\\0&1\end{pmatrix}\) are diagonalizable.\end{exercise}
\begin{exercise}[Further, Challenging] Prove that an operator is
diagonalizable exactly when the sum of the dimensions of its eigenspaces is
\(\dim V\).\end{exercise}
\begin{exercise}[Core] For the standard basis of \(F^3\), compute its dual
basis and the matrix of the dual of a specified \(3\times3\) matrix.  Then
verify the congruence formula for the bilinear form whose matrix is
\[
\begin{pmatrix}2&1\\1&0\end{pmatrix}
\]
after the basis change \((e_1,e_2)\mapsto(e_1+e_2,e_2)\).
\end{exercise}
